\documentclass[12pt,a4paper]{article}
\usepackage[utf8]{inputenc}
\usepackage[indonesian]{babel}
\usepackage{amsmath}
\usepackage{amsfonts}
\usepackage{amssymb}
\usepackage{geometry}
\usepackage{hyperref}

\geometry{margin=2.5cm}

\title{Metode Newton-Raphson Modifikasi dalam Komputasi Numerik}
\author{Ahmad Fakhrul Bawani}
\date{}

\begin{document}

\maketitle

\section{Pendahuluan}
Metode Newton-Raphson standar memiliki kelemahan utama ketika berhadapan dengan akar ganda (akar dengan multiplisitas $m > 1$). Pada kondisi ini, konvergensi yang seharusnya kuadratik akan melambat menjadi linear. Metode Newton-Raphson Modifikasi dirancang untuk mengembalikan kecepatan konvergensi kuadratik tersebut.

\section{Konsep dan Rumus}
Terdapat dua pendekatan utama dalam modifikasi metode ini:

\subsection{Pendekatan 1: Jika Multiplisitas ($m$) Diketahui}
Jika kita mengetahui bahwa akar tersebut memiliki multiplisitas $m$ (misalnya akar ganda dua, maka $m=2$), rumus iterasinya menjadi:
\begin{equation}
x_{n+1} = x_n - m \frac{f(x_n)}{f'(x_n)}
\end{equation}

\subsection{Pendekatan 2: Menggunakan Fungsi Bantu $u(x)$}
Jika multiplisitas $m$ tidak diketahui, kita dapat mendefinisikan fungsi baru $u(x)$:
\begin{equation}
u(x) = \frac{f(x)}{f'(x)}
\end{equation}
Akar dari $u(x)$ adalah sama dengan akar dari $f(x)$, namun akar ganda pada $f(x)$ akan menjadi akar tunggal pada $u(x)$. Dengan menerapkan metode Newton-Raphson pada $u(x)$, kita mendapatkan rumus:
\begin{equation}
x_{n+1} = x_n - \frac{f(x_n)f'(x_n)}{[f'(x_n)]^2 - f(x_n)f''(x_n)}
\end{equation}

\section{Kelebihan dan Kekurangan}
\subsection{Kelebihan}
\begin{itemize}
  \item Sangat efektif untuk menyelesaikan persamaan yang memiliki akar ganda.
  \item Mengembalikan laju konvergensi kuadratik yang hilang pada metode standar.
\end{itemize}

\subsection{Kekurangan}
\begin{itemize}
  \item Memerlukan perhitungan turunan kedua $f''(x)$ untuk pendekatan kedua, yang secara komputasi bisa sangat mahal atau sulit diturunkan.
  \item Lebih kompleks dalam implementasi kode dibandingkan metode standar.
\end{itemize}

\end{document}